\chapter{Experiments and results}\label{AER}

\section{Algorithms and Heuristics}

\subsection{Exhaustive}
The simplest strategy for finding the best MCM for a dataset, is to try every combination of ICC's possible for the amount of variables in the dataset. This is as of now the only guaranteed way to find the globally optimal MCM for a dataset. While this naive algorithm works well enough on small systems, it scales poorly with the amount of variables in the system. Each variable added roughly doubles the amount of MCMs which need to be checked. The algorithm therefore scales as \(O(2^n)\), with \(n\) the amount of variables in the dataset. The algorithm is mostly an adaptation of the code by~\citet{bayesianinference}.

\subsection{Greedy Algorithms}
When the modeled systems become too large, the requirement of finding te globally optimal solution needs to be loosened. This allows for the use of algorithms such as the greedy algorithms described in this subsection. A greedy algorithm can be defined as a process in which a solution gets improved in steps. The algorithm can perform a certain operation each step. All of the states a solution can reach within a single step are called the neighborhood of a solution. At each step, the state in the neighborhood which improves on the current solution is chosen and passed on to the next step. This is repeated until the solution cannot be improved anymore.

A good example of a greedy algorithm is what we call the merge algorithm and is also called hierarchical merging in the paper by~\citet{bayesianinference}. In this algorithm, the MCM starts as system in which each variable is in its own ICC. These ICCs are then greedily merged together. To determine the best merge in the neighborhood, the algorithm needs to consider every combination of two ICCs. This is \(\frac{1}{2}n^2 - n\) comparisons in the worst case, and it has to run at the worst \(n\) steps, which gives this algorithm a complexity of \(O(n^3)\).

Another example is what we call the constructive algorithm. This algorithm constructs solutions from an empty MCM. Variables are then added one by one, greedily picking whether to add the variable to one of the existing ICCs or to create a new one. The variables can be added in a fixed order, or the variable to add can be picked greedily. With a fixed order, the algorithm needs to consider at worst \(n\) ICCs for \(n\) steps, which gives this algorithm a complexity of \(O(n^2)\). The greedy order variant also needs to consider at worst \(n\) variables per step, so it has a complexity of \(O(n^3)\).
% honestly not sure about these complexities they may be nonsense

Lastly, there is the local search algorithm, though this procedure could also be classified as a heuristic. The local search algorithm can start with an arbitrary MCM, either a trivial one or a result from one of the previously explained algorithms. That MCM then gets improved by moving a single variable to a different or new ICC each step, continuing until no more improvements can be found. This does mean that this the amount of steps that this algorithm will take before reaching an optimum cannot be predicted beforehand. This means that an algorithmic complexity cannot be given for this algorithm/heuristic.

The greedy behavior of the above algorithms often leads them to end up in local minima. This raises the question whether the algorithm would have avoided these local minima if it had taken a less optimal choice earlier in the process. To bring insight into this question, all of the greedy algorithms have the option to keep a list of the best found MCMs in a step. The neighborhoods of all of these MCMs is then searched to find a list of candidates for the next step. The maximum size of this list is called the beam size, and can be set by the user. This concept is borrowed from beam search, a lossy optimization of the exhaustive breadth-first search, in which only the best candidates in a step are kept and the rest discarded.

\subsection{Metaheuristics}
When algorithms or heuristics on their own cannot deliver a satisfactory result to a problem in a timely manner, it can be beneficial to employ the use of metaheuristics. Metaheuristics are processes with can guide or augment an underlying algorithm or heuristic to achieve better results. An example of this in our library is tabu search. This metaheuristic modifies the local search algorithm by adding a ``tabu list''. This is a list of ICCs which new candidate solutions cannot have. ICCs get added to this list, when they are left by the solution. In practice this means that, when a new solution gets accepted where a variable moves from one ICC to another, the old states of these ICCs gets added to the tabu list. The length of the tabu list is limited, meaning that previously banned ICCs become available again after a number of steps. This number will likely be slightly higher than the half of the tabu list length, as almost all modifications end up changing two ICCs at a time. The last modification to the local search algorithm is that tabu search is allowed to accept moves which make the held solution worse than before, if these moves are the only ones allowed by the tabu list. This has the consequence that it can force the algorithm to walk out of local minima and explore the surrounding solution space.

The following metaheuristics all use a strategy of making a random mutation to an existing solution and then accepting or rejecting this random mutation based on some criteria. The mutation performed on the existing solution is generated the same way for all of the following metaheuristics. The generated mutation can be one of three possible kinds: a split, merge, or swap. The swap is almost identical to the moves performed by the local search explained in the previous section; a random variable is moved to a different or new ICC. The merge mutation merges two random ICCs together, similarly to the operation performed by the merge search. The split mutation selects a random amount of variables from an ICC and splits them off into a new ICC. These three categories are generally chosen with equal probability, except in the edge cases of a trivial or one-ICC system, in which case merge or split is chosen, respectively.

The simplest heuristic to build with this mutation principle is the hill climbing heuristic. This heuristic uses the accepting criteria that a mutation can only be accepted if the solution after the mutation has a better Log(E) than before. This causes the solution to climb up in solution space, until it reaches a local minimum. Because the mutations are generated randomly, the heuristic will often converge on different minima every time it is run, but it lacks the ability to walk out of these minima to explore the surrounding solution space. For better exploration, it needs to be augmented so that the solution can accept temporarily worse solutions.

Late acceptance hill climbing is a metaheuristic which tries to implement this augmentation.